\documentclass[11pt,a4paper]{article}

\usepackage[margin=1.05in]{geometry}
\usepackage{amsmath,amssymb,amsthm}
\usepackage{mathtools}
\usepackage{booktabs}
\usepackage{array}
\newcolumntype{L}[1]{>{\raggedright\arraybackslash}p{#1}}
\usepackage{xcolor}
\usepackage[colorlinks=true,linkcolor=black,citecolor=black,urlcolor=blue]{hyperref}
\usepackage{enumitem}
\usepackage[T1]{fontenc}
\usepackage[protrusion=true,expansion=false]{microtype}
\setlength{\emergencystretch}{2.5em}

\theoremstyle{plain}
\newtheorem{theorem}{Theorem}[section]
\newtheorem{proposition}[theorem]{Proposition}
\newtheorem{lemma}[theorem]{Lemma}
\newtheorem{conjecture}[theorem]{Conjecture}
\theoremstyle{definition}
\newtheorem{definition}[theorem]{Definition}
\newtheorem{example}[theorem]{Example}
\newtheorem{question}[theorem]{Open question}
\theoremstyle{remark}
\newtheorem{remark}[theorem]{Remark}

%%%%%%%%%%%%%%%%%%%%%%%%%%%%%%%%%%%%%%%%%%%%%%%%%%%%%%%%%%%%%%%%%%%%%%%%%%%%%%%
% NOTATION RULES OBSERVED THROUGHOUT:
%   * the rho calculus is never written with the Greek letter rho
%   * quoting is @P, dropping is *x; corner quotes are pre-2016 notation
%%%%%%%%%%%%%%%%%%%%%%%%%%%%%%%%%%%%%%%%%%%%%%%%%%%%%%%%%%%%%%%%%%%%%%%%%%%%%%%
\newcommand{\rhoc}{\textsf{rho}}
\newcommand{\pr}[2]{\langle #1,\,#2\rangle}
\newcommand{\qn}[2]{@\langle #1,\,#2\rangle}
\newcommand{\out}[2]{#1!(#2)}
\newcommand{\rcv}[3]{\mathsf{for}(#1 \leftarrow #2)\,#3}
\newcommand{\drop}[1]{{*}#1}
\newcommand{\bdrop}[2]{{*}(#1,\,#2)}
\newcommand{\unq}[1]{\mathop{\&}#1}
\newcommand{\esub}[2]{[\,#1/#2\,]}
\newcommand{\scope}[2]{(#1)\,#2}
\newcommand{\red}{\rightsquigarrow}
\newcommand{\scong}{\equiv_{S}}
\newcommand{\ssub}[2]{\{#1/#2\}}
\newcommand{\ysub}[2]{\{#1/#2\}_{\mathrm{syn}}}
\newcommand{\Sem}{\textsf{Sem}[D]}
\newcommand{\Unq}{\textsf{Unq}[D]}
\newcommand{\ES}{\textsf{ES}[D]}
\newcommand{\fv}{\mathrm{fv}}
\newcommand{\barb}[1]{\mathop{\downarrow}_{#1}}
\newcommand{\nbarb}[1]{\mathop{\not\downarrow}_{#1}}

\title{Where Substitution Happens\\[0.4em]
\large Semantic substitution, deferred unquote and explicit substitution\\
in a paired-message \rhoc{} calculus}
\author{L.G. Meredith\\ F1R3FLY.io}
\date{Working note, 1 October 2026}

\begin{document}
\maketitle

\begin{abstract}
The note \emph{Unitary in the Equations, Measurement in the Reductions}
\cite{UE} arrived, on the way to a quantum interpretation, at a small
extension of the \rhoc{} calculus that stands on its own: messages and names
are pairs of a process and a datum drawn from a type $D$, dropping splits into
a binary equation that preserves the datum and a unary drop resolved by
semantic substitution, and replication is derived by reflection. This note
extracts that core, independently of any quantum reading, and uses it to
compare three treatments of information exchange at interaction. In the first,
the one adopted, the received code is installed by semantic substitution at
the communication itself. In the second, substitution is syntactic and a
separate unary rule $\unq{\qn{p}{d}}\red p$ unquotes later. In the third, the
substitution is a process term left beside the continuation, in the style of
explicit fusions. The two alternatives were proposed by Mike Stay. The
argument against them has three parts: unary reductions that no interaction
triggers do not fit the context-labelled bisimulation programme; both
alternatives change the shape of bisimulation; and that change is the formal
shadow of a practical cost, because once substitution is detached from
communication the programmer must coordinate by hand which substitutions
happen together and which happen in order. Worked examples show both burdens:
on the programmer, and on the bisimulation proofs. A generic catalyst
transformation, which turns any unary rule $A\red B$ into $A\mid C\red B$,
answers the first objection; the note shows that it leaves the other two
standing, and that it makes the remaining coordination burden explicit rather
than removing it. The conclusion ties the argument to the rho combinators:
the size of translations from languages with variables into combinator
languages is usually read as evidence of the compression variables offer, but
when delivery of values is concurrent part of that compression is a
consolidation of coordination, and semantic substitution is what achieves it.
\end{abstract}

\tableofcontents

%=============================================================================
\section{Introduction}
%=============================================================================

Every process calculus must say what happens to the information carried by a
message at the moment it is received. In the \rhoc{} calculus \cite{MR05} the
question has an extra edge, because what is received is a quoted process, and
the receiver may do two quite different things with it: treat it as a name, or
run it. The original paper answers with \emph{semantic substitution}:
substitution is syntactic everywhere except at a drop $\drop{y}$, where it
installs the received process in place. The quantum design session recorded in
\cite{UE} rediscovered this choice under pressure, and in doing so produced a
version of the calculus in which the choice is unusually easy to see. Names
and messages are pairs, so there is a datum to lose; the binary drop keeps it
and the unary drop erases it; and nothing else in the calculus touches the
question.

That makes the core calculus a good place to compare alternatives. In
conversation, Mike Stay proposed two. The first keeps substitution uniformly
syntactic and moves execution of received code into a separate unary
reduction. The second makes substitution a process term, after Wischik and
Gardner's explicit fusions \cite{GW00}. Both are natural, and both are standard
moves elsewhere in the literature on substitution \cite{ACCL91}. This note
records why the paired-message calculus keeps semantic substitution, and makes
the reasons concrete with examples.

\paragraph{Reading order.} Section~\ref{sec:core} states the core calculus.
Section~\ref{sec:alts} states the two alternatives. Section~\ref{sec:where}
compares all three by asking where the unquote, and with it the erasure of the
datum, takes place. Section~\ref{sec:argument} states the argument.
Sections~\ref{sec:programmer} and~\ref{sec:bisim} give the examples, first the
programmer's burden and then the burden on bisimulation proofs.
Section~\ref{sec:fair} records what the alternatives do offer.
Section~\ref{sec:conclusion} draws the conclusion, connecting the argument
to the rho combinators, and Section~\ref{sec:open} lists what remains open.

%=============================================================================
\section{The core calculus}\label{sec:core}
%=============================================================================

The construction is a type parameterisation of the \rhoc{} calculus. Fix a
nonempty set $D$ of data, and a distinguished $d_0\in D$. Nothing below uses
any structure on $D$; in \cite{UE} it is the group $\mathbb{Z}/8$ of phase
labels, but here it is arbitrary.

\begin{definition}[Syntax of $\Sem$]\label{def:syntax}
\[
\begin{array}{rcll}
M & ::= & \pr{P}{d} & \text{messages},\ d\in D\\[2pt]
x,u & ::= & @M \mid y & \text{names}\\[2pt]
P & ::= & 0 \mid \out{x}{M} \mid \rcv{y}{x}{P} \mid P\mid P
      \mid \bdrop{x}{u} \mid \drop{y} & \text{processes}
\end{array}
\]
The receipt $\rcv{y}{x}{P}$ binds $y$ in $P$. The unary drop $\drop{y}$
applies only to variables. There is no restriction operator and no persistent
receipt: fresh names arise by quotation, and replication is derived
(Proposition~\ref{prop:repl}).
\end{definition}

\begin{definition}[Structural congruence]\label{def:cong}
$\scong$ is the least congruence, closed under $\alpha$-conversion, satisfying
\begin{description}[leftmargin=2.5em,style=nextline]
\item[(E1) Monoid] $P\mid 0\scong P$,\quad $P\mid Q\scong Q\mid P$,\quad
  $(P\mid Q)\mid R\scong P\mid(Q\mid R)$.
\item[(E2) Binary drop] $\bdrop{\qn{p}{d}}{u}\scong \out{u}{\pr{p}{d}}$.
\end{description}
Names are identified up to the congruence of what they quote:
$\qn{p}{d}=\qn{q}{e}$ iff $p\scong q$ and $d=e$.
\end{definition}

\begin{remark}[Why E2 is an equation]
E2 could have been read as a notational abbreviation instead. It is folded into
$\scong$ for two reasons. As an equation it is a bijection between term
shapes: every output $\out{u}{M}$ is congruent to exactly one binary drop
$\bdrop{@M}{u}$, and nothing is lost in either direction. And because it sits in
the quotient, the observation theory of Section~\ref{sec:bisim}, which is
defined up to $\scong$, never has to distinguish a forwarded message from its
expansion. The binary drop forwards the whole pair; in particular it
\emph{preserves the datum}.
\end{remark}

\begin{definition}[Semantic substitution]\label{def:sub}
Let $N=\qn{p}{d}$. The substitution $\ssub{N}{y}$ acts on names by
$y\ssub{N}{y}=N$, $z\ssub{N}{y}=z$ for $z\neq y$, and
$\qn{Q}{e}\ssub{N}{y}=\qn{Q\ssub{N}{y}}{e}$. On processes it commutes with
every constructor, renaming bound variables to avoid capture, with one
exception:
\[
  (\drop{y})\ssub{\qn{p}{d}}{y} \;=\; p .
\]
In particular $\bdrop{y}{u}\ssub{N}{y}=\bdrop{N}{u}\scong\out{u}{\pr{p}{d}}$
keeps the datum, while the unary drop discards it. Substitution acts under
quotation, so a name built from received code changes at the moment the code
arrives.
\end{definition}

\begin{definition}[Reduction]\label{def:red}
\[
\begin{array}{ll}
\textsc{(Comm)} & \rcv{y}{x}{P}\mid\out{x}{\pr{q}{d}}
                  \;\red\; P\ssub{\qn{q}{d}}{y}\\[4pt]
\textsc{(Par)} & P\red P' \;\Rightarrow\; P\mid Q\red P'\mid Q\\[4pt]
\textsc{(Struct)} & P\scong P',\ P'\red Q',\ Q'\scong Q
                    \;\Rightarrow\; P\red Q
\end{array}
\]
COMM is the only rule that does any work, and every reduction is the meeting of
a sender and a receiver.
\end{definition}

\begin{proposition}[Closure under reflection]
Every communication substitutes a name of the form $@M$, and E2 takes such a
name back to a message. The payload sort and the quoted sort coincide.
\end{proposition}

\begin{proposition}[Replication is derivable]\label{prop:repl}
Let $R_x := \rcv{y}{x}{(\bdrop{y}{x}\mid\drop{y})}$ and
$!P := \out{x}{\pr{R_x\mid P}{d_0}}\mid R_x$. Then
\[
  !P \;\red\; \out{x}{\pr{R_x\mid P}{d_0}}\mid R_x\mid P .
\]
\end{proposition}

\begin{proof}
COMM substitutes $N=\qn{R_x\mid P}{d_0}$ for $y$. The binary drop becomes
$\bdrop{N}{x}\scong\out{x}{\pr{R_x\mid P}{d_0}}$, which re-emits the message
with its datum intact, and the unary drop becomes $R_x\mid P$.
\end{proof}

The proof shows the two drops doing different jobs within a single step: one
forwards a copy for the next unfolding, the other installs a copy to run. Both
happen at once, because both are effects of the same substitution.

\begin{remark}[Installation and execution]\label{rem:install}
Semantic substitution fixes the moment received code is \emph{installed}: at
the COMM. It does not force the code to run then. If $\drop{y}$ is guarded,
say $\rcv{y}{x}{\rcv{z}{c}{\drop{y}}}$, then $p$ is installed under the guard
at the first COMM and runs when the guard is lifted by the second. Both
moments are communications, and both are under the program's control.
\end{remark}

%=============================================================================
\section{The two alternatives}\label{sec:alts}
%=============================================================================

\subsection{Deferred unquote}

\begin{definition}[$\Unq$]
Replace the unary drop with an unquote operator, and substitute
syntactically:
\[
\begin{array}{rcl}
P & ::= & 0 \mid \out{x}{M} \mid \rcv{y}{x}{P} \mid P\mid P
      \mid \bdrop{x}{u} \mid \unq{x}\\[4pt]
\textsc{(Comm$_{\textsf{U}}$)} & & \rcv{y}{x}{P}\mid\out{x}{\pr{q}{d}}
      \;\red\; P\ysub{\qn{q}{d}}{y}\\[4pt]
\textsc{(Unq)} & & \unq{\qn{p}{d}} \;\red\; p
\end{array}
\]
with E1, E2, PAR and STRUCT as before.
\end{definition}

The operator must apply to arbitrary names, not only to variables, because
syntactic substitution turns $\unq{y}$ into $\unq{\qn{q}{d}}$. A consequence is
that $\unq{\qn{p}{d}}$ can be written directly in source text, as a term that
reduces with nobody else involved. In $\Sem$ the corresponding term is never
formed. The binary drop is unaffected: substitution into $\bdrop{y}{u}$ is
already syntactic, and E2 applies.

\subsection{Explicit substitution}

The second proposal leaves the substitution itself in the term, more or less
following explicit fusions \cite{GW00,PV98}. The rules below are a
reconstruction, made to have something definite to argue about; the precise
rules intended may differ, and the argument is written so that it does not
depend on the details.

\begin{definition}[$\ES$]
Extend the syntax of $\Unq$ with a substitution term and a scope for the
variable it governs:
\[
\begin{array}{rcl}
P & ::= & \cdots \mid \esub{N}{y} \mid \scope{y}{P}\\[4pt]
\textsc{(Comm$_{\textsf{E}}$)} & & \rcv{y}{x}{P}\mid\out{x}{M}
      \;\red\; \scope{y}{(P\mid\esub{@M}{y})}
\end{array}
\]
with UNQ for unquoting, and the fusion-style equations
\begin{description}[leftmargin=2.5em,style=nextline]
\item[(X1) Propagation] $\scope{y}{(P\mid\esub{N}{y})}\scong
  \scope{y}{(P\ysub{N}{y}\mid\esub{N}{y})}$.
\item[(X2) Collection] $\scope{y}{\esub{N}{y}}\scong 0$.
\item[(X3) Extrusion] $(\scope{y}{P})\mid Q\scong\scope{y}{(P\mid Q)}$ when
  $y\notin\fv(Q)$.
\end{description}
\end{definition}

The scope is needed. In the receipt, $y$ is bound by the input prefix; COMM$_{\textsf{E}}$ removes the prefix, so something else must keep $y$ from
becoming a free, global variable. In the fusion calculus that job is done by
restriction, which keeps fusions local. The \rhoc{} calculus has no
restriction, deliberately, since fresh names arise by quotation, so $\ES$ must
import one.

There are two ways to read the propagation of the substitution term, and they
lead to different calculi.
\begin{itemize}[leftmargin=1.5em]
\item \textbf{Equational reading.} X1 holds as stated. Then by X1, X3 and X2,
  $\scope{y}{(P\mid\esub{N}{y})}\scong P\ysub{N}{y}$, so $\ES$ coincides with
  $\Unq$ up to $\scong$ and adds only the bookkeeping of the scope. Everything
  said about $\Unq$ below then applies to it unchanged.
\item \textbf{Interactional reading.} X1 is dropped, and each unguarded
  occurrence of $y$ (as a channel subject, or under $\unq{}$) is resolved by its
  own reduction against the substitution term, which persists until collected:
  for instance $\esub{\qn{p}{d}}{y}\mid\unq{y}\red\esub{\qn{p}{d}}{y}\mid p$.
  Each resolution is now an interaction, which answers the first objection
  below, but a single COMM is followed by one resolution step per occurrence,
  and occurrences under quotation are never reached, because nothing reduces
  inside a name.
\end{itemize}

%=============================================================================
\section{Where the unquote lives}\label{sec:where}
%=============================================================================

The three calculi agree on everything except the treatment of the unary drop.
The binary drop is the same equation in all of them. Without the unary drop,
nothing runs received code, recursion is lost, and with linear receipts every
term strongly normalises \cite{UE}. So the unary drop is exactly where
recursion comes from, and the whole comparison concerns one operator.

The unary drop is also the one place the datum is erased:
$\qn{p}{d}\mapsto p$ is many-to-one. It is therefore useful to ask, of each
design, which step performs that erasure. A fourth placement, the equation
$\drop{\qn{p}{d}}\scong p$, completes the picture.

\begin{center}
\begin{tabular}{@{}L{0.30\textwidth}L{0.16\textwidth}L{0.16\textwidth}L{0.28\textwidth}@{}}
\toprule
\textbf{Placement} & \textbf{Triggered by interaction} & \textbf{Atomic with the COMM} & \textbf{Where $d$ is erased}\\
\midrule
Equation, $\drop{\qn{p}{d}}\scong p$ & no step & no step & in $\scong$, at no cost\\
Semantic substitution ($\Sem$) & yes & yes & at the COMM\\
Deferred unquote ($\Unq$) & no & no & at a UNQ step, per occurrence\\
Deferred unquote, catalysed ($\Unq^{C}$) & yes, with the catalyst & no & at a catalysed step, per occurrence\\
Explicit substitution ($\ES$), equational & no & no & as in $\Unq$\\
Explicit substitution ($\ES$), interactional & yes & no & at each resolution\\
\bottomrule
\end{tabular}
\end{center}

The equational placement is ruled out in \cite{UE} because equations are free
while information erasure is not: an equation that forgets $d$ makes $\scong$
identify configurations that differ in data, and cost accounting has nowhere to
charge for the loss. Semantic substitution charges the erasure to the COMM,
which is already the step that interaction pays for. The deferred unquote
charges it to a step that no interaction pays for, unless it is catalysed
(Section~\ref{sec:catalyst}), in which case the catalyst pays. Explicit
substitution either inherits that placement or spreads the erasure across as
many steps as there are occurrences.

%=============================================================================
\section{The argument}\label{sec:argument}
%=============================================================================

\begin{enumerate}[leftmargin=1.5em]
\item \textbf{Reductions should be interactions.} A unary rule like
  $\unq{\qn{p}{d}}\red p$, which fires without any interaction triggering it,
  does not fit the context-labelled bisimulation programme \cite{LM00}. There,
  a transition is labelled by the minimal context that completes a redex, and
  every redex of $\Sem$ is a sender and receiver meeting. A redex that is
  complete on its own contributes transitions that no context takes part in,
  and that the environment can neither cause nor see except through their
  effects. This objection, unlike the other two, has a generic answer, the
  catalyst transformation of Section~\ref{sec:catalyst}.
\item \textbf{Both proposals change bisimulation.} $\Unq$ adds internal steps
  between a communication and the running of what it delivered. $\ES$ puts
  substitution terms into the states being compared. Either way the
  equivalence, and the proofs about it, must be rethought.
\item \textbf{That change reflects a cost to the programmer.} In $\Sem$ the
  substitution happens atomically, at the COMM, and the programmer
  coordinates when and where substitutions happen with the full power of the
  language, because every substitution is a communication. In the
  alternatives, substitutions are detached from the communications that
  produced them, and the programmer has to do extra work to make sure that
  substitutions that must coincide do, and substitutions that must be ordered
  are.
\end{enumerate}

Points 2 and 3 are one point seen from two sides: the more complicated
bisimulation is the formal record of the coordination the programmer now has
to do. The next two sections give examples of each side.

%=============================================================================
\section{The burden on the programmer}\label{sec:programmer}
%=============================================================================

Throughout, $a$, $b$, $c$, $x$, $x_1$, $x_2$ and $\mathit{ack}$ are fixed,
distinct quoted names, and $e\in D$ is an arbitrary datum. Write $P\barb{a}$ when
$P\scong\out{a}{M}\mid Q$ for some $M$ and $Q$, so that $P$ offers an output
on $a$.

\subsection{Knowing that received code is live}

\begin{example}[Install, then acknowledge]\label{ex:ack}
A receiver installs code it is sent and then acknowledges:
\[
  S \;=\; \out{x}{\pr{p}{d}}\mid
         \rcv{y}{x}{(\drop{y}\mid\out{\mathit{ack}}{\pr{0}{e}})}.
\]
In $\Sem$,
$S\red p\mid\out{\mathit{ack}}{\pr{0}{e}}$. Whoever consumes the
acknowledgement does so in a state where $p$ is already installed: the
acknowledgement means what it says. In $\Unq$,
\[
  S \;\red\; \unq{\qn{p}{d}}\mid\out{\mathit{ack}}{\pr{0}{e}},
\]
and the acknowledgement can be consumed, and acted on, before UNQ fires.
\end{example}

The natural repair is for the receiver to acknowledge \emph{after} the
unquote. That cannot be done from the receiver's side. Suppose the receiver
wraps the code, sends the wrapper to itself, and unquotes it:
\[
  \rcv{y}{x}{\big(\out{w}{\pr{\unq{y}\mid\out{\mathit{ack}}{\pr{0}{e}}}{e}}
     \mid\rcv{z}{w}{\unq{z}}\big)}.
\]
After both communications and the UNQ of the wrapper, the state is
$\unq{\qn{p}{d}}\mid\out{\mathit{ack}}{\pr{0}{e}}$, which is where it started.
Every wrapper leaves the inner unquote pending.

\begin{proposition}[No receiver-side signal of liveness]\label{prop:noack}
In $\Unq$, no part of the receiver can be causally downstream of the step
$\unq{\qn{p}{d}}\red p$, except through actions of $p$ itself.
\end{proposition}

\begin{proof}[Proof sketch]
UNQ has a single residual, $p$, and no partner. Anything else in the term is
unchanged by the step and so cannot depend on it. The only events that follow
it causally are those of $p$.
\end{proof}

So in $\Unq$ the acknowledgement must be built into the shipped code: the
\emph{sender} must send $\pr{p\mid\out{\mathit{ack}}{\pr{0}{e}}}{d}$. That moves
a receiver's coordination problem across the trust boundary to the sender,
and it changes the message, hence the name $@M$, hence any protocol keyed on
that name. In $\Sem$ the receiver needs no help. Its continuation running is
the evidence that the code is installed, because both are the same step.

\subsection{Substitutions that must happen in order}

\begin{example}[Activation follows reception]\label{ex:order}
A receiver takes two pieces of code in sequence and runs each:
\[
  T \;=\; \out{x}{\pr{p_1}{d_1}}\mid\out{x}{\pr{p_2}{d_2}}\mid
       \rcv{y_1}{x}{\big(\drop{y_1}\mid\rcv{y_2}{x}{\drop{y_2}}\big)}.
\]
Taking $p_1$ as the first one received, in $\Sem$ every run passes through
$p_1\mid\rcv{y_2}{x}{\drop{y_2}}\mid\out{x}{\pr{p_2}{d_2}}$: the first code
is installed before the second can be received. In $\Unq$ this run is
available:
\[
\begin{array}{rcl}
T & \red & \unq{\qn{p_1}{d_1}}\mid\rcv{y_2}{x}{\unq{y_2}}\mid\out{x}{\pr{p_2}{d_2}}\\
  & \red & \unq{\qn{p_1}{d_1}}\mid\unq{\qn{p_2}{d_2}}\\
  & \red & \unq{\qn{p_1}{d_1}}\mid p_2\\
  & \red & p_1\mid p_2 .
\end{array}
\]
The second code is live while the first is not.
\end{example}

The order of reception no longer fixes the order of activation. When that
order matters, as when $p_2$ is an update meant to supersede $p_1$, the
programmer must restore it, and by Proposition~\ref{prop:noack} the only way
is the one from Example~\ref{ex:ack}: the sender of $p_1$ builds in a
signal, and the receiver guards the second receipt on it.

\subsection{Substitutions that must happen together}

\begin{example}[Simultaneous availability]\label{ex:together}
Two pieces of code arrive on different channels, from different senders, and
should become available together. Let $p_1=\out{a}{\pr{0}{e}}$ and
$p_2=\out{b}{\pr{0}{e}}$, and
\[
  U \;=\; \out{x_1}{\pr{p_1}{d_1}}\mid\out{x_2}{\pr{p_2}{d_2}}\mid
       \rcv{y_1}{x_1}{\rcv{y_2}{x_2}{(\drop{y_1}\mid\drop{y_2})}}.
\]
In $\Sem$, $p_1$ is installed under the guard on $x_2$ at the first COMM, and
both are released by the second: $U\red\red p_1\mid p_2$. Along every run,
a state offers output on $a$ only if it also offers output on $b$:
$P\barb{a}$ implies $P\barb{b}$. In $\Unq$, after both COMMs and one UNQ the
state is $p_1\mid\unq{\qn{p_2}{d_2}}$, with $\barb{a}$ but $\nbarb{b}$.
\end{example}

As in the previous examples, the receiver cannot repair this. Packaging the
two received codes into one wrapper, $\out{w}{\pr{\unq{y_1}\mid\unq{y_2}}{e}}$,
and unquoting it leaves $\unq{\qn{p_1}{d_1}}\mid\unq{\qn{p_2}{d_2}}$, two
independent pending steps again. In $\Unq$, two codes received from
independent senders cannot be made live simultaneously by the receiver alone.
In $\Sem$ it takes no effort at all, because a single COMM installs every
drop of the variables it has bound so far.

With $k$ unguarded drops of received variables, a single COMM of $\Sem$
corresponds in $\Unq$ to a COMM followed by $k$ independent unquotes, with
$2^k$ reachable states of partial activation. Each of these is a state the
programmer must either show to be harmless or rule out.

\subsection{Names computed from received code}

Because substitution in $\Sem$ acts under quotation, a name built from
received code is resolved at the COMM along with everything else. This is how
\rhoc{} programs compare code: two processes rendezvous on a name exactly when
the code they quote is congruent.

\begin{example}[Naming the code that arrived]\label{ex:name}
Let $n(y)=\qn{\drop{y}}{d_0}$, a name computed from received code, and
\[
  V \;=\; \out{x}{\pr{p}{d}}\mid\rcv{y}{x}{\out{n(y)}{\pr{0}{e}}}
       \mid\rcv{z}{\qn{p}{d_0}}{Q}.
\]
In $\Sem$, COMM gives $n(y)\ssub{\qn{p}{d}}{y}=\qn{p}{d_0}$, so the output
meets the third component and $Q$ runs. In $\Unq$ the output lands on
$\qn{\unq{\qn{p}{d}}}{d_0}$, which is a different name from $\qn{p}{d_0}$,
because UNQ is a reduction, not an equation, and nothing reduces inside a
quote. The rendezvous never happens.
\end{example}

The workaround in $\Unq$ is to listen on $\qn{\unq{\qn{p}{d}}}{d_0}$. But that
name depends on the datum $d$ of the message in which the code happened to
travel. The same code obtained by three routes, received with datum $d$,
received with datum $d'$, or written literally, gets three different names.
Code identity becomes route-dependent, and the programmer must track the route
of every piece of code whose name is used.

$\ES$ under the interactional reading adds a variant of this burden even for
channel positions. A standard \rhoc{} idiom is to reply on the name of the
request itself:
\[
  \out{x}{M}\mid\rcv{r}{@M}{S}\;\mid\;\rcv{y}{x}{\out{y}{\pr{q}{e}}}.
\]
In $\Sem$ the server's reply is on $@M$ as soon as the request is received.
Under the interactional reading the reply sits on the variable $y$ until a
separate resolution step against $\esub{@M}{y}$, and if the server replies
twice, the two resolutions are independent and the simultaneity problem of
Example~\ref{ex:together} returns.

%=============================================================================
\section{The burden on bisimulation}\label{sec:bisim}
%=============================================================================

\subsection{Unary redexes and labels}

In the approach of Leifer and Milner \cite{LM00}, an agent $a$ has a
transition $a\xrightarrow{C}a'$ when $C[a]\red a'$ and $C$ is a minimal
context that makes this happen; their congruence theorem is for the strong
equivalences built on these labels. In $\Sem$ every redex consists of a
receiver and a sender, so a label is either the missing half of a
communication or the identity, when both halves are already present in the
agent. In $\Unq$ the term $\unq{\qn{p}{d}}$ is a complete redex on its own. It
contributes identity-labelled transitions that correspond to no communication
at all.

Strong equivalences count these steps. $\unq{\qn{p}{d}}$ has exactly one
transition and its target is $p$, so it is not strongly bisimilar to $p$
except in degenerate cases. Every law of $\Sem$ that involves running
received code therefore survives in $\Unq$ at best as a weak law, and
extending the context-labelled congruence results from strong to weak
equivalences requires further conditions that must then be checked.

\subsection{The catalyst transformation}\label{sec:catalyst}

Mike Stay pointed out that the first objection has a generic answer. It comes
from a device the author introduced in 2014 to repair a categorical semantics
for the pi-calculus, published as the catalysis construction of \cite{SM15},
where it restricts the 2-morphisms interpreting rewrites to the contexts in
which they may apply. The idea is to make a unary reduction depend on a \emph{catalyst}, a fixed
process whose presence enables it.

\begin{definition}[Catalyst transformation]
Fix a process $C$. Every rule of the form $A\red B$ in which no interaction
triggers the step is replaced by
\[
  A\mid C\;\red\;B .
\]
In the reactive system, the step then has the labelling context
$[\,]\mid C$. In the \emph{catalytic} variant the catalyst is returned,
$A\mid C\red B\mid C$.
\end{definition}

Applied to $\Unq$, this gives a calculus $\Unq^{C}$ with
\[
  \textsc{(Unq$_C$)}\qquad \unq{\qn{p}{d}}\mid C\;\red\;p
  \qquad\text{or}\qquad
  \unq{\qn{p}{d}}\mid C\;\red\;p\mid C .
\]
Now $\unq{\qn{p}{d}}$ is no longer a complete redex on its own. Its transition
is
\[
  \unq{\qn{p}{d}}\xrightarrow{\;[\,]\mid C\;}p ,
\]
labelled by the context that supplies the catalyst, and every reduction of $\Unq^{C}$ is an interaction. The
first objection is met in full, and generically: any calculus with unary rules
can be brought inside the context-labelled programme this way.

The interactional reading of $\ES$ is an instance of the catalytic variant,
with the substitution term itself as the catalyst:
$\esub{N}{y}\mid\unq{y}\red\esub{N}{y}\mid p$. This explains why that reading
already answered the first objection in Section~\ref{sec:alts}.

The transformation makes each unquote an interaction. It does not make the
unquote part of the communication that delivered the code, and the remaining
objections are about exactly that.
\begin{itemize}[leftmargin=1.5em]
\item \textbf{Atomicity is not restored.} A COMM followed by $k$ unguarded
  unquotes is still $k+1$ steps, now $k$ of them catalysed. The $2^k$ partial
  states of Example~\ref{ex:together} and of the next subsection remain, and so
  does the simultaneity problem: one catalysed step enables one unquote.
\item \textbf{Quotation is untouched.} No reduction acts inside a quote,
  catalysed or not, so Example~\ref{ex:name} and
  Proposition~\ref{prop:namefail} hold for $\Unq^{C}$ as stated.
\item \textbf{Strong equivalences still separate $\unq{\qn{p}{d}}$ from $p$.}
  The step is now labelled $[\,]\mid C$ instead of the identity, but it is
  still a step, so strong context-labelled bisimilarity still distinguishes
  the pending unquote from the running code.
\end{itemize}

What changes is who can control the timing. If $C$ is ambient, for instance a
catalyst that is always present and always returned, the labels change but the
programmer gains no handle on when unquotes fire. If $C$ is something the
program emits and withholds, say an output $\out{\kappa}{\pr{0}{d_0}}$ on a
name $\kappa$ that the program holds, then the program can decide when each
unquote is allowed to happen. That helps with ordering, but only partly.
To activate $p_2$ only after $p_1$ in Example~\ref{ex:order}, the receiver
needs evidence that the first unquote has completed. Under $A\mid C\red B$
there is none: the catalyst is consumed and nothing new appears. Under the
catalytic variant the catalyst reappears, but its presence does not
distinguish ``not yet used'' from ``used and returned'', and a receiver that
waits on it competes with the unquote for it. So
Proposition~\ref{prop:noack} carries over. Evidence of completion exists only
if the catalysed rule is designed to emit a signal, $A\mid C\red B\mid C'$.

A catalyst that signals completion is, however, a small synchronisation
protocol laid alongside every unquote: release the catalyst, wait for the
signal, proceed. That is the coordination of the third objection, written out
explicitly. The catalyst transformation therefore does not remove the
programmer's burden; it makes the burden visible as protocol. It shows that
what separates the alternatives from semantic substitution is not the absence
of an interaction but the absence of atomicity. $\Sem$ needs no catalyst,
because its only unquote is already part of the one interaction that
delivered the code.

\subsection{The relations get larger}

Consider a receiver with $k$ unguarded drops of the received variable,
\[
  W \;=\; \out{x}{\pr{p}{d}}\mid\rcv{y}{x}{(\drop{y}\mid\cdots\mid\drop{y})}.
\]
In $\Sem$, $W$ has one reduction, to $p^{\mid k}$, and any bisimulation
containing $W$ needs one more pair to account for it. In $\Unq$, $W$ reduces
to $(\unq{\qn{p}{d}})^{\mid k}$, and then through every subset of the $k$
unquotes. A weak bisimulation relating $W$ to its $\Sem$ counterpart must
contain a pair for each of the $2^k$ partial states, each related to
$p^{\mid k}$ or to an intermediate of $p$'s own behaviour. These are the same
$2^k$ states that burden the programmer in Example~\ref{ex:together}, now
appearing as proof obligations. Up-to techniques can compress them, but
they have to be developed and justified, whereas in $\Sem$ they are not
needed.

\subsection{Laws that fail outright}

The extra steps of $\Unq$ can be absorbed by moving to weak equivalence. The
name problem of Example~\ref{ex:name} cannot.

\begin{proposition}\label{prop:namefail}
Let
\[
  S_1 = \out{x}{\pr{p}{d}}\mid\rcv{y}{x}{\out{\qn{\drop{y}}{d_0}}{\pr{0}{e}}},
  \]
  \[
  S_2 = \out{x}{\pr{p}{d}}\mid\rcv{y}{x}{\out{\qn{p}{d_0}}{\pr{0}{e}}},
\]
with $y\notin\fv(p)$, and $\drop{y}$ read as $\unq{y}$ in $\Unq$. Then $S_1$
and $S_2$ are strongly barbed bisimilar in $\Sem$, and not even weakly barbed
bisimilar in $\Unq$.
\end{proposition}

\begin{proof}[Proof sketch]
In $\Sem$ both have the barb $\barb{x}$ and a single reduction, to the same
term $\out{\qn{p}{d_0}}{\pr{0}{e}}$. In $\Unq$ the reducts are
$\out{\qn{\unq{\qn{p}{d}}}{d_0}}{\pr{0}{e}}$ and
$\out{\qn{p}{d_0}}{\pr{0}{e}}$. Both are stuck, and their barbs are on
distinct names, since $\unq{\qn{p}{d}}\not\scong p$. No sequence of internal
steps reconciles them, because no step acts inside a quote.
\end{proof}

The bisimilarity in $\Sem$ is not a congruence: a context can send different
code on $x$. The point is that in $\Unq$ the two terms are distinguished even
without a context. The law ``receiving code and naming it is the same as
naming the code'' holds in $\Sem$ and fails in $\Unq$, however weak the
equivalence. Bisimulation proofs in $\Unq$ must keep track of the route by
which each named piece of code arrived, which is Example~\ref{ex:name}
restated as a proof obligation.

\subsection{Substitutions in the state}

Under the equational reading, $\ES$ reduces to $\Unq$ up to $\scong$, so
bisimulation up to $\scong$ absorbs the substitution terms and the previous
subsections apply unchanged. Under the interactional reading the substitution
terms are genuine components of the state. A bisimulation must relate
$\scope{y}{(P\mid\esub{N}{y})}$, and each of its partial resolutions, to
$P\ssub{N}{y}$, and the number of internal steps a single COMM gives rise to
depends on the number of occurrences of $y$, which is a syntactic accident.
Names with a free variable, such as $\qn{\unq{y}}{d_0}$, mean different
channels under different surrounding substitutions, so barbs must be read
relative to their scope. This is why the fusion calculus works with
hyperequivalence, bisimulation closed under substitutions \cite{PV98}. That is
a heavier notion than anything $\Sem$ needs, since in $\Sem$ no state ever
contains an unresolved received name.

\subsection{Cost}

In cost-accounted \rhoc{}, cost is charged to reductions. In $\Sem$ the cost
of a run is a function of its communications. In $\Unq$ it also depends on how
many drops the receiver happened to write, and under the interactional reading
of $\ES$ on how many occurrences of the variable there are in any position. A
cost-aware equivalence therefore distinguishes programs that differ only in how
often they mention what they received. This is the cost-accounting version of
the remark in Section~\ref{sec:where}: the alternatives charge for erasure at
steps that interaction does not pay for.

%=============================================================================
\section{What the alternatives offer}\label{sec:fair}
%=============================================================================

The alternatives have real merits, and the argument does not depend on
denying them.

\begin{itemize}[leftmargin=1.5em]
\item \textbf{Provenance.} In $\Unq$ a name built from received code records
  the datum of the message the code arrived in. Sometimes that is wanted. In
  $\Sem$ it is available on request rather than imposed: the binary drop keeps
  the datum, so $\qn{\bdrop{y}{u}}{d_0}$, which is congruent to
  $\qn{\out{u}{\pr{p}{d}}}{d_0}$ after the COMM, is a name that records
  $d$. The programmer chooses whether a name records where the code came from.
\item \textbf{Sharing.} Explicit substitution earns its place in
  implementations, where it avoids copying large terms and lets a substitution
  be shared \cite{ACCL91}. An implementation of $\Sem$ may well use explicit
  substitutions internally. The argument here is about the calculus, which
  fixes what is observable: whatever the machinery, the installation of
  received code must look atomic with the communication that delivered it.
\item \textbf{Uniform substitution.} $\Unq$ has a simpler substitution
  operation, which is pleasant metatheoretically. The price is paid
  elsewhere: in an extra reduction rule, in weak equivalences, and in the
  failure of Proposition~\ref{prop:namefail}.
\end{itemize}

%=============================================================================
\section{Conclusion: two kinds of compression}\label{sec:conclusion}
%=============================================================================

Translating a term language with variables and nested scopes into a
combinator language, which is much flatter, typically blows up the size of
terms. The effect is especially striking for Yoshida's concurrent combinators
for the pi-calculus \cite{Yoshida02}. On the family
$P_d=\rcv{x_1}{a}{\rcv{x_2}{a}{\cdots\rcv{x_d}{a}{0}}}$, of size $d+1$, her
prefix elimination produces $(3^{d-1}+1)/2$ atoms \cite{RC3}. It is natural
to take such blow-ups as strong evidence for the compression that variables
provide.

The usual account of that compression is abstraction. A variable is a single
point of reference, and a binder consolidates the mechanism that carries a
value from the environment to its place in the term. SKI has no such
mechanism \cite{CF58}. Bracket abstraction rebuilds the route from the outside
of a term to each use of a value on the inside, again and again: at every
application node, and once more for every level of nesting. Measured against
SKI, the value of variables appears to be the consolidation of these
\emph{routes}.

There is a richer perspective, and this note is what exposes it. When the
calculations that carry a value to its uses are allowed to be concurrent,
there is a second cost beyond routing: \emph{coordination}. Delivering a value
to several occurrences becomes several independent events. Someone must then
arrange that deliveries which have to coincide do, and deliveries which have
to be ordered are. Sections~\ref{sec:programmer} and~\ref{sec:bisim} show
that this cost is real, both for the programmer and in the proofs, and the
catalyst transformation of Section~\ref{sec:catalyst} shows that paying it
amounts to laying a small protocol alongside every occurrence. Semantic
substitution consolidates these coordination costs. A single interaction
delivers the value to every occurrence at once, and the coordination is borne
by the communication that was going to happen anyway. This is a richer form of
compression than abstraction alone: variables with semantic substitution
compress not only the description of \emph{where} a value goes, but also the
protocol by which it arrives \emph{everywhere at once}.

Yoshida's blow-up can be read in these terms. In the recurrence
$T(j+1)=3T(j)-1$ behind the count above, each level of nesting adds two kinds
of atom to a body of $n$ atoms. Rule~(I) adds $n-1$ duplicators, which split
the body at every parallel composition: that is routing, the part of the cost
SKI would recognise. Rule~(VI) puts every atom behind a synchroniser of its
own, adding $n$ more: that is coordination, one synchronisation per
occurrence, which is exactly the per-occurrence structure of the deferred
unquote. Composed down a chain of nested prefixes, both compound. The
reflective encoding of \cite{RC3} reaches $3d+1$ by quoting a continuation
once and releasing it through a three-atom gate. The gate is a single
synchronisation for the whole continuation, which is the combinator-level
counterpart of installing received code by semantic substitution.
Reflection does for the combinators what semantic substitution does for the
term language: it consolidates coordination.

The comparison in this note also isolates the coordination component cleanly.
$\Sem$, $\Unq$ and $\ES$ all keep variables and binders, so routing is the
same in each. Whatever a faithful translation from $\Sem$ into one of the
alternatives costs, it is the price of coordination alone, with abstraction
held fixed.

%=============================================================================
\section{Open questions}\label{sec:open}
%=============================================================================

\begin{question}[Operational correspondence]
Translate $\Sem$ into $\Unq$ by reading $\drop{y}$ as $\unq{y}$.
Proposition~\ref{prop:namefail} shows the translation is not adequate in
general. Is it adequate up to weak bisimilarity on the fragment in which no
drop occurs under quotation? The examples in Section~\ref{sec:programmer}
suggest it is, but this has not been proved.
\end{question}

\begin{question}[Relative pushouts with E2]
With E2 in $\scong$, every output is congruent to a binary drop. Labels can
then be taken in a normal form with binary drops on quoted names expanded.
It remains to check that the reactive system for $\Sem$ has the relative
pushouts that the congruence theorem of \cite{LM00} requires, and to state
how semantic substitution interacts with the contexts involved.
\end{question}

\begin{question}[The explicit-substitution rules]
The rules for $\ES$ here are a reconstruction. With the intended rules in
hand, it should be checked which reading they fall under, and whether a
middle reading exists that both keeps every resolution an interaction and
restores atomicity.
\end{question}

\begin{question}[Which catalyst]
For $\Unq^{C}$ the choice of catalyst matters: ambient or emitted, consumed or
returned, silent or signalling. With a signalling catalyst emitted by the
receiver, is there a systematic translation of $\Sem$ into $\Unq^{C}$ that
reconstructs atomicity by protocol, and what does it cost, in steps and in the
size of the bisimulations needed to prove it correct?
\end{question}

\begin{question}[Measuring the two compressions]
Can the split of Section~\ref{sec:conclusion} be made precise, as two
measures on translations: one counting routing (duplicators, path
construction) and one counting coordination (synchronisers, catalysed
steps)? A natural first test is the translation of $\Sem$ into $\Unq^{C}$
with a signalling catalyst, where routing is unchanged and every cost of the
translation should register as coordination; a second is to apportion the
$\Theta(3^{d})$ of Yoshida's elimination between the two measures.
\end{question}

\begin{question}[Structure on $D$]
Nothing here uses structure on $D$. The quantum note \cite{UE} uses a group;
other choices (costs, provenance tags, types) suggest themselves. The question
is which properties of the core are stable under which choices, in particular
whether erasure at the unary drop should be parametric in a map $D\to 1$ that
the choice of $D$ can refine.
\end{question}

\paragraph{Acknowledgement.} The two alternatives, and the conversation that
sharpened the argument, are due to Mike Stay. The note was drafted in dialogue
with Claude (Anthropic).

\begin{thebibliography}{99}
\bibitem{ACCL91} M.~Abadi, L.~Cardelli, P.-L.~Curien and J.-J.~L\'evy.
  Explicit substitutions. \emph{J. Funct. Program.} 1(4):375--416, 1991.
\bibitem{CF58} H.~B.~Curry and R.~Feys. \emph{Combinatory Logic}, vol.~I.
  North-Holland, 1958.
\bibitem{GW00} P.~Gardner and L.~Wischik. Explicit fusions. \emph{MFCS 2000},
  LNCS 1893, 2000.
\bibitem{LM00} J.~J.~Leifer and R.~Milner. Deriving bisimulation congruences
  for reactive systems. \emph{CONCUR 2000}, LNCS 1877, 2000.
\bibitem{UE} L.~G.~Meredith. Unitary in the equations, measurement in the
  reductions. Working note, F1R3FLY-io/publications, \texttt{fuzzyware/}, 2026.
\bibitem{MR05} L.~G.~Meredith and M.~Radestock. A reflective higher-order
  calculus. \emph{ENTCS} 141(5):49--67, 2005.
\bibitem{PV98} J.~Parrow and B.~Victor. The fusion calculus: expressiveness and
  symmetry in mobile processes. \emph{LICS 1998}.
\bibitem{RC3} F1R3FLY.io. Name-free combinators for the rho calculus.
  Research note, draft~3, F1R3FLY-io/publications, \texttt{rho-combinators/},
  2026.
\bibitem{SM15} M.~Stay and L.~G.~Meredith. Higher category models of the
  pi-calculus. \emph{Higher-Dimensional Rewriting and Applications (HDRA)},
  Warsaw, 2015. arXiv:1504.04311.
\bibitem{Yoshida02} N.~Yoshida. Minimality and separation results on
  asynchronous mobile processes: representability theorems by concurrent
  combinators. \emph{Theoret. Comput. Sci.} 274(1--2):231--276, 2002.
\end{thebibliography}

\end{document}
